%%%%%%%%%%%%%%%%%%%%%%%%%%%%%%%%%%%%%%%%%%%%%%%%%%%%%%%%%%%%%%%%%%%%%%%%%%%
% Handout und Uebungsblatt zu Lehreinheit 8
% Datensatz (Aufgabe 2) ist fiktiv; Kennwerte im Erwartungshorizont mit
% Python (numpy/scipy) nachgerechnet.
%%%%%%%%%%%%%%%%%%%%%%%%%%%%%%%%%%%%%%%%%%%%%%%%%%%%%%%%%%%%%%%%%%%%%%%%%%%
\begin{handout}{Quantitative Methoden und Statistik}

\textbf{Worum es heute ging.} Drei Becher, eine Kugel: Die Spielleitung
öffnet wissentlich einen leeren Becher. Wechseln gewinnt in zwei von drei
Fällen, obwohl \enquote{zwei Becher, also fifty-fifty} sich zwingend
anfühlt. Öffnet die Spielleitung dagegen blind, ist es tatsächlich egal.
Dasselbe Bild, zwei Prozesse, zwei Wahrscheinlichkeiten. Statistik ist
das Werkzeug, mit dem wir solche Täuschungen der Intuition prüfen.

\textbf{Kernbegriffe}
\begin{description}
  \item[Bedingte Wahrscheinlichkeit] $P(A\mid B)$: Wahrscheinlichkeit
    von $A$, wenn $B$ bekannt ist. Die Richtung zählt:
    $P(\text{Daten}\mid H)\neq P(H\mid\text{Daten})$.
  \item[Deskriptiv / Inferenz] Daten beschreiben bzw.\ von der Stichprobe
    auf die Grundgesamtheit schließen.
  \item[p-Wert] Wahrscheinlichkeit für mindestens so extreme Daten,
    \emph{falls} die Nullhypothese stimmt. Kein Maß für die Wahrheit
    einer Hypothese und keins für die Größe eines Effekts
    \parencite{wassersteinlazar2016asa}.
  \item[Effektgröße] Größe eines Unterschieds unabhängig von $n$, z.\,B.\
    Cohens $d$ (0,2 klein, 0,5 mittel, 0,8 groß).
  \item[Power] Wahrscheinlichkeit, einen vorhandenen Effekt zu finden.
    Für $d=0{,}5$ und Power 0,80 braucht ein $t$-Test 64 Personen pro
    Gruppe.
\end{description}

\textbf{Welcher Kennwert passt?}

\begin{center}
\small
\begin{tabular}{@{}llll@{}}
\toprule
\textbf{Skala} & \textbf{Beispiel} & \textbf{Lage} & \textbf{Streuung} \\
\midrule
nominal & Strategie, Vorwissen ja/nein & Modus & Häufigkeiten \\
ordinal & Staunen 1--5 & Median & Quartile, $\mathit{IQR}$ \\
metrisch & Ausschlag in mm & Mittelwert, Median & $s$, $\mathit{IQR}$, Spannweite \\
\bottomrule
\end{tabular}
\end{center}
Ausreißerregel nach \textcite{tukey1977eda}: Werte, die mehr als
$1{,}5\cdot\mathit{IQR}$ unter $Q_1$ oder über $Q_3$ liegen.

\begin{merkbox}
Erst fragen, wie die Daten zustande kamen, dann rechnen. Ein kleiner
p-Wert ist kein Beweis, und unter vielen Tests ist er fast garantiert.
\end{merkbox}

\textbf{Riskante Vorhersagen.} Folgt eine Beobachtung $E$ aus der
Hypothese $H$, dann gilt $P(H\mid E)/P(H)=1/P(E)$: Je unwahrscheinlicher
die Vorhersage war, desto mehr beweist ihr Eintreffen
\parencite{horwich1982probability}.

\textbf{Zum Weiterlesen:} \textcite{doeringbortz2016}, Kap.~12;
\textcite{wassersteinlazar2016asa}; zum Ziegenproblem
\textcite{mhall_krauss2003}.
\end{handout}

\begin{uebungsblatt}{Becher, Daten, p-Werte}

\aufgabe{Monty Hall in Paaren}[in der Stunde, Paare, 10 Min.]
Die Spielleitung würfelt verdeckt die Lage der Kugel (1--2: A, 3--4: B,
5--6: C), die spielende Person zeigt auf einen Becher, die Spielleitung
öffnet einen \emph{leeren}, nicht gewählten Becher. Ungerade Runden:
bleiben, gerade Runden: wechseln. Kreuzen Sie an, ob gewonnen wurde.

\begin{center}
\small
\renewcommand{\arraystretch}{1.25}
\begin{tabular}{|c|c|c|c||c|c|c|c|}
\hline
\textbf{Rd.} & \textbf{Strategie} & \textbf{Kugel} & \textbf{gewonnen} &
\textbf{Rd.} & \textbf{Strategie} & \textbf{Kugel} & \textbf{gewonnen} \\ \hline
1 & bleiben & & & 13 & bleiben & & \\ \hline
2 & wechseln & & & 14 & wechseln & & \\ \hline
3 & bleiben & & & 15 & bleiben & & \\ \hline
4 & wechseln & & & 16 & wechseln & & \\ \hline
5 & bleiben & & & 17 & bleiben & & \\ \hline
6 & wechseln & & & 18 & wechseln & & \\ \hline
7 & bleiben & & & 19 & bleiben & & \\ \hline
8 & wechseln & & & 20 & wechseln & & \\ \hline
9 & bleiben & & & 21 & bleiben & & \\ \hline
10 & wechseln & & & 22 & wechseln & & \\ \hline
11 & bleiben & & & 23 & bleiben & & \\ \hline
12 & wechseln & & & 24 & wechseln & & \\ \hline
\end{tabular}
\end{center}
Summe: \textbf{bleiben} \rule{12mm}{0.3pt} von \rule{12mm}{0.3pt}
gewonnen \qquad \textbf{wechseln} \rule{12mm}{0.3pt} von
\rule{12mm}{0.3pt} gewonnen \qquad Kurs gesamt (Tafel):
\rule{25mm}{0.3pt}

\aufgabe{Datensatz-Analyse}[in der Stunde und Hausaufgabe, Kleingruppe]
\textbf{Fiktiver Datensatz} (konstruiert für diese Übung, keine echte
Erhebung): 50 Personen haben ein Ideomotor-Pendel gehalten (vgl.\ LE~1).
Gemessen wurde der größte Ausschlag in Millimetern (\textbf{mm}), danach
gaben die Personen an, wie stark sie gestaunt haben (\textbf{St}, 1 = gar
nicht bis 5 = sehr), und ob sie den Pendeleffekt vorher kannten
(\textbf{VW}, 1 = ja, 0 = nein).

\begin{center}
\footnotesize
\setlength{\tabcolsep}{3.2pt}
\begin{tabular}{@{}rrrr@{\hspace{9pt}}rrrr@{\hspace{9pt}}rrrr@{\hspace{9pt}}rrrr@{\hspace{9pt}}rrrr@{}}
\toprule
\textbf{Nr} & \textbf{mm} & \textbf{St} & \textbf{VW} &
\textbf{Nr} & \textbf{mm} & \textbf{St} & \textbf{VW} &
\textbf{Nr} & \textbf{mm} & \textbf{St} & \textbf{VW} &
\textbf{Nr} & \textbf{mm} & \textbf{St} & \textbf{VW} &
\textbf{Nr} & \textbf{mm} & \textbf{St} & \textbf{VW} \\
\midrule
 1 & 11 & 3 & 0 & 11 &  9 & 3 & 0 & 21 & 10 & 2 & 1 & 31 & 43 & 3 & 1 & 41 & 47 & 4 & 0 \\
 2 &  5 & 2 & 0 & 12 & 24 & 2 & 0 & 22 & 26 & 4 & 1 & 32 & 57 & 4 & 0 & 42 & 59 & 3 & 0 \\
 3 & 15 & 3 & 1 & 13 & 24 & 4 & 1 & 23 & 39 & 4 & 0 & 33 & 21 & 3 & 0 & 43 & 26 & 3 & 0 \\
 4 & 41 & 4 & 1 & 14 &  8 & 2 & 0 & 24 & 15 & 4 & 0 & 34 & 96 & 5 & 0 & 44 & 31 & 4 & 1 \\
 5 & 37 & 3 & 0 & 15 & 16 & 2 & 1 & 25 & 13 & 3 & 0 & 35 & 49 & 4 & 0 & 45 & 23 & 2 & 0 \\
 6 & 39 & 5 & 0 & 16 & 25 & 4 & 0 & 26 & 13 & 3 & 0 & 36 & 30 & 4 & 0 & 46 & 20 & 3 & 1 \\
 7 & 18 & 3 & 0 & 17 & 19 & 5 & 0 & 27 & 66 & 5 & 0 & 37 & 23 & 4 & 0 & 47 & 12 & 1 & 1 \\
 8 &  0 & 1 & 0 & 18 & 53 & 4 & 1 & 28 & 35 & 3 & 1 & 38 & 53 & 4 & 0 & 48 & 36 & 4 & 1 \\
 9 & 16 & 2 & 1 & 19 &  5 & 2 & 1 & 29 & 38 & 5 & 0 & 39 & 18 & 2 & 0 & 49 & 18 & 3 & 0 \\
10 & 31 & 4 & 1 & 20 & 45 & 5 & 0 & 30 & 26 & 3 & 1 & 40 & 73 & 5 & 0 & 50 & 40 & 4 & 0 \\
\bottomrule
\end{tabular}
\end{center}

\begin{enumerate}
  \item Bestimmen Sie für jede Variable das Skalenniveau.
  \item Berechnen Sie für \textbf{mm}: Minimum, Maximum, Spannweite,
    Median, $Q_1$, $Q_3$, $\mathit{IQR}$ und Mittelwert. Geben Sie an,
    wie Sie die Quartile bestimmt haben.
  \item Zeichnen Sie ein Histogramm (Klassenbreite 10\,mm) und einen
    Boxplot. Gibt es Ausreißer nach der $1{,}5\cdot\mathit{IQR}$-Regel?
    Was würden Sie mit ihnen tun?
  \item Berechnen Sie Mittelwert und Median von \textbf{mm} ohne den
    Ausreißer. Welcher Kennwert reagiert stärker, und warum?
  \item Fassen Sie \textbf{St} sinnvoll zusammen. Ist ein Mittelwert
    zulässig?
  \item Vergleichen Sie \textbf{mm} zwischen Personen mit und ohne
    Vorwissen mit Mittelwert \emph{und} Median. Was fällt auf?
\end{enumerate}
\schreiblinien{6}

\aufgabe{Eine Zeitungsmeldung lesen}[in der Stunde, Plenum]
\textbf{Fiktive Meldung} (konstruiert für diese Übung):

\begin{quote}\small
\emph{Zaubertricks machen klüger -- jetzt bewiesen.} Wer vor dem Lernen
einen Zaubertrick sieht, behält mehr. Das zeigt eine Online-Befragung
unter 112 Studierenden. Die Forschenden erhoben 23 Merkmale des
Lernverhaltens. Bei der Frage \enquote{Ich konnte mich gut
konzentrieren} lag die Gruppe mit Zaubertrick signifikant vorn ($p =
0{,}04$). \enquote{Damit ist der Nutzen von Magie im Hörsaal
wissenschaftlich belegt}, sagt die Studienleiterin.
\end{quote}
Beantworten Sie: Was wurde gemessen, an wem? Was wird behauptet? Wird
eine Effektgröße genannt? Wie viele Vergleiche wurden gerechnet, und was
folgt daraus für den einen signifikanten Befund?
\schreiblinien{4}

\aufgabe{Methodische Kurzeinschätzung}[Hausaufgabe, max.\ 10 Sätze]
Finden Sie in einer Tageszeitung oder einem Nachrichtenportal einen
Artikel, der eine Studie mit $p<0{,}05$ als \enquote{Beweis} darstellt.
Schreiben Sie eine methodische Kurzeinschätzung mit Quellenangabe. Nutzen
Sie die vier Fragen aus Aufgabe~3.

\aufgabe{Tagebuch des Staunens}[Hausaufgabe]
Wann hat Ihr Gefühl für Wahrscheinlichkeit Sie zuletzt getäuscht?
Notieren Sie Erwartung, Ereignis und den passenden Begriff aus dieser
Einheit.

\begin{lehrhinweis}[Erwartungshorizont]
\textbf{Aufgabe~1.} Erwartet sind pro Paar etwa 8 von 12 Gewinnen beim
Wechseln und 4 von 12 beim Bleiben. Einzelne Paare streuen stark (die
Standardabweichung der Gewinnzahl bei 12 Wechselrunden beträgt etwa 1,6),
manche Paare finden sogar das Gegenteil. Erst der Kurs-Pool zeigt das
Muster stabil. Das ist der Anlass, über Stichprobengröße zu sprechen.

\textbf{Aufgabe~2} (Werte mit Python nachgerechnet).
(1)~Nr: nominal (Kennung); mm: Verhältnisskala; St: ordinal; VW: nominal
(dichotom).
(2)~$n=50$; Minimum 0, Maximum 96, Spannweite 96\,mm; Median 25,5\,mm;
Mittelwert 29,94\,mm; Standardabweichung $s=19{,}32$\,mm. Quartile als
Median der unteren bzw.\ oberen Hälfte: $Q_1=16$, $Q_3=40$,
$\mathit{IQR}=24$. Mit linearer Interpolation (z.\,B.\ Excel
\texttt{QUARTILE.INKL}) ergibt sich $Q_3=39{,}75$ und
$\mathit{IQR}=23{,}75$; beides ist richtig, wenn das Verfahren genannt
wird. Modus: 18 und 26 (je dreimal).
(3)~Klassenhäufigkeiten 0--9: 5; 10--19: 13; 20--29: 10; 30--39: 9;
40--49: 6; 50--59: 4; 60--69: 1; 70--79: 1; 80--89: 0; 90--99: 1. Die
Verteilung ist rechtsschief. Zäune: $16-36=-20$ und $40+36=76$. Einziger
Ausreißer ist Nr.~34 mit 96\,mm; 73\,mm (Nr.~40) liegt knapp innerhalb.
Mit Interpolation liegen die Zäune bei $-19{,}6$ und $75{,}4$; Ergebnis
gleich. Gute Antworten löschen den Wert nicht stillschweigend, sondern
prüfen ihn (Messfehler? jemand hat bewusst geschwungen?) und berichten
beide Auswertungen.
(4)~Ohne Nr.~34: Mittelwert 28,59\,mm, Median 25,0\,mm, $s=16{,}98$\,mm.
Der Mittelwert sinkt um 1,35\,mm, der Median um 0,5\,mm, die
Standardabweichung um 2,3\,mm.
(5)~Häufigkeiten St: 1: 2; 2: 9; 3: 15; 4: 17; 5: 7. Median 3, Modus 4.
Ein Mittelwert (3,36) ist streng genommen nicht zulässig, weil die
Abstände der Stufen unbekannt sind; wer ihn angibt, sollte das begründen.
Optional: Spearmans Rangkorrelation zwischen mm und St beträgt
$\rho=0{,}71$.
(6)~Mit Vorwissen ($n=17$): Mittelwert 25,88\,mm, Median 26\,mm,
$s=13{,}06$. Ohne Vorwissen ($n=33$): Mittelwert 32,03\,mm, Median
25\,mm, $s=21{,}75$. Der Mittelwert legt einen Unterschied von gut 6\,mm
nahe, der Median fast keinen; die Gruppe ohne Vorwissen streut stärker
und enthält die beiden größten Werte (73 und 96\,mm). Die erwartete
Antwort: Aus diesen Daten folgt kein belastbarer Gruppenunterschied, und
die Wahl des Kennwerts bestimmt die Geschichte, die man erzählen könnte.

\textbf{Aufgabe~3.} Gemessen: Selbstauskunft, keine Lernleistung;
Online-Befragung, unklar, ob die Gruppen zufällig gebildet wurden. 23
Merkmale: Bei $\alpha=0{,}05$ erwartet man unter reinem Zufall gut einen
\enquote{signifikanten} Befund ($23\cdot0{,}05\approx1{,}2$); die
Wahrscheinlichkeit für mindestens einen liegt bei $1-0{,}95^{23}\approx
69\,\%$. Keine Effektgröße. \enquote{Bewiesen} ist nichts; allenfalls
liegt eine Hypothese für eine präregistrierte Folgestudie vor (Brücke zu
LE~10).
\end{lehrhinweis}
\end{uebungsblatt}
